\documentclass[10pt,twocolumn]{article}
\usepackage[letterpaper,top=0.8in,bottom=0.85in,left=0.75in,right=0.75in,columnsep=0.25in]{geometry}
\usepackage{lmodern}
\usepackage{microtype}
\usepackage{amsmath,amssymb}
\usepackage{booktabs}
\usepackage{enumitem}
\usepackage{tabularx}
\usepackage{listings}
\usepackage{xcolor}
\usepackage[bookmarks=true,bookmarksopen=true,bookmarksnumbered=true,colorlinks=true,linkcolor=blue!70!black,citecolor=green!50!black,urlcolor=blue!80!black]{hyperref}

\pagestyle{plain}

\title{\textbf{\Large TerraShift+: A Physics-Informed Machine Learning and Multi-Mission NASA Earth Observation Pipeline for Climate-Resilient Crop Rotation Optimization}}

\author{
  \textbf{Lazaro A. Philemon}$^{1}$ \quad \textbf{Jackson Haruni}$^{1}$ \quad \textbf{Arnold Philemon}$^{2}$ \\
  $^{1}$\textit{Tanzania} \quad $^{2}$\textit{Kenya} \\
  \textit{Jiajiri Organization --- Agronomy, Remote Sensing \& Systems Engineering Collective} \\
  \href{https://github.com/jiajiri/TerraShift}{\texttt{https://github.com/jiajiri/TerraShift}}
}

\date{\small NASA Space Apps Challenge 2026 Submission $\bullet$ Challenge: Field Shift: Adapting Farms with NASA Data $\bullet$ November 2026}

\begin{document}

\maketitle

\begin{abstract}
\noindent
Smallholder and commercial agricultural production systems worldwide confront compound climate stressors: erratic monsoonal onsets, prolonged mid-season dry spells, extreme temperature surges, accelerated topsoil detachment, and catastrophic soil organic carbon (SOC) depletion. Traditional decision support systems rely upon static historical crop almanacs that fail under non-stationary climatic extremes, or alternatively, ungrounded black-box Large Language Models (LLMs) prone to agronomic hallucinations and mass-balance violations. In this paper, we present \textbf{TerraShift+}, an audited physics-informed machine learning (PIML) and closed-loop decision support engine that ingests ten distinct NASA Earth Observation (EO) satellite missions and global pedotransfer datasets. TerraShift+ couples empirical satellite telemetry with classical biophysical governing laws: FAO-56 Dual $K_c$ Penman-Monteith reference evapotranspiration ($ET_0$), Saxton-Rawls pedotransfer hydrology for available water capacity ($AWC$), dynamic root-zone water balance and Crop Water Stress Index ($WSI$), Growing Degree Days ($GDD$) and Killing Degree Days ($KDD$), the Revised Universal Soil Loss Equation (RUSLE), RothC-modeled carbon balance ($\Delta SOC$), and Log-Logistic Standardized Precipitation-Evapotranspiration Index (SPEI). Surrogate yield predictions are delivered via spatial-cross-validated Quantile Gradient Boosting Regressors ($P10/P50/P90$), guarded by a Gaussian Mixture Model (GMM) Mahalanobis Out-of-Distribution (OOD) distance metric. Optimal multi-year crop sequences are resolved across four distinct operational archetypes (Commercial Maximum Profit, Climate Adaptation \& Drought-Resilient, Soil Restorative \& Conservation, and Balanced Sustainable) using a multi-objective Pareto solver evaluating 100-draw Monte Carlo weather shocks and 10th-percentile Conditional Value at Risk ($\text{CVaR}_{10}$). Over 72 spatial clusters, the surrogate attains an RMSE of 0.050, an $R^2$ of 0.942 against FAO-56 Penman-Monteith, pinball losses of 0.009 ($P10$) and 0.008 ($P90$), and empirical 80\% coverage of 83.5\%. Stress-testing across canonical control stations (Dodoma, Tanzania; Fallon, Nevada; Ames, Iowa) and a polar negative control confirms strict monotonicity under drought shock ($-25\%$ precipitation, $+2.0^\circ\text{C}$ warming) and reliable OOD rejection ($D_M = 8.40$). Server-Sent Event (SSE) streaming delivers full 5-year multi-objective optimizations in 1.50~ms (cached) and 585~ms (uncached cold start), providing smallholder farmers and extension officers with physically bounded, climate-resilient agronomic guidance.
\end{abstract}

\noindent\textbf{Keywords:} crop rotation, NASA POWER, SMAP L3/L4, GPM IMERG, MODIS, ECOSTRESS, FAO-56 Penman-Monteith, Saxton-Rawls pedotransfer, RUSLE, quantile regression, CVaR, Out-of-Distribution detection, physics-informed machine learning, agricultural decision support systems.

\section{Introduction}

\subsection{Background on Climate Volatility \& Bare Soil Vulnerability}
Global agriculture is entering an era of unprecedented climate non-stationarity. The thermodynamic intensification of atmospheric circulation regimes has drastically amplified the frequency, duration, and severity of compound agroclimatic extremes. Smallholder and commercial agricultural producers across sub-Saharan tropical savannas, Mediterranean climatic basins, and temperate continental corn belts no longer operate within the stationary precipitation and temperature envelopes established during the late twentieth century. In place of predictable seasonal onsets, growing cycles are routinely disrupted by delayed monsoonal rains, punctuated by high-intensity convective storms that shear topsoil aggregates, and terminated by unseasonal dry spells coupled with extreme vapor pressure deficits.

A primary accelerator of agricultural collapse during these erratic climate shocks is the widespread persistence of continuous monoculture cropping systems paired with prolonged bare-soil fallow windows. When arable land is left devoid of vegetative canopy cover between harvesting and subsequent planting, the soil surface is directly subjected to the kinetic energy of unattenuated raindrop impact. Raindrop shear destroys structural soil macro-aggregates, dispersing clay and silt particles that wash into surface pores to create an impermeable physical crust. This crusting phenomenon suppresses hydraulic infiltration rates, converting precious precipitation into destructive surface runoff.

Under the physics of the Revised Universal Soil Loss Equation (RUSLE), the lack of surface vegetative cover ($C$-factor approaching unity) exposes fertile topsoil to kinetic erosion rates exceeding 10 to 25 metric tons per hectare annually on modest slope gradients. Concurrently, unbroken monocropping of nitrogen-exhaustive cereals (e.g., continuous maize or grain sorghum) rapidly depletes root-zone organic matter and macro-nutrients, necessitating escalating applications of costly synthetic fertilizers. For smallholder households operating on narrow margins, this creates a vulnerability cycle of soil depletion and crop failure: rising input costs paired with climate-induced harvest failures precipitate acute rural debt cycles and chronic food insecurity.

\subsection{Closed-Loop System Overview \& Architecture}
To overcome the structural failures of static historical almanacs and unconstrained generative AI models, we introduce \textbf{TerraShift+}, an audited, closed-loop decision support engine that operationalizes planetary-scale NASA Earth observations through physics-informed machine learning (PIML).

The architecture couples a Python 3.14 computational core with a React 18 + Vite visual cockpit deployed serverlessly on Cloudflare Pages. In place of opaque deep learning pipelines, TerraShift+ embeds physical conservation laws directly into its feature transformations, surrogate regression models, and sequence optimization algorithms:

\begin{enumerate}[leftmargin=*]
  \item \textbf{NASA Multi-Mission Ingestion Pipeline:} Programmatically assimilates near-real-time and multi-decadal observations across ten remote sensing and pedological assets, including NASA POWER v2.0 daily agroclimatology, SMAP L4 root-zone soil wetness, ECOSTRESS thermal Land Surface Temperature (LST), GPM IMERG calibrated precipitation, MODIS MOD13Q1 vegetative indices, GEDI L4A aboveground biomass, and ISRIC/iSDAsoil 30m--250m profiles.
  \item \textbf{Audited Biophysical Agronomy Engine:} Replaces empirical heuristics with explicit governing equations: FAO-56 Dual $K_c$ Penman-Monteith evapotranspiration, Saxton-Rawls pedotransfer hydrology, finite-depth root-zone bucket water balance, phenological Growing/Killing Degree Days, RUSLE soil detachment, and RothC turnover kinetics.
  \item \textbf{Quantile Regression ML Surrogate:} Employs spatial-cross-validated Quantile Gradient Boosting Regressors to predict 5-year yields across calibrated uncertainty bounds ($P10, P50, P90$), providing risk-averse farmers with transparent downside protections.
  \item \textbf{GMM Mahalanobis Out-of-Distribution Guardrail:} Evaluates input feature vectors against a 72-cluster Gaussian Mixture Model to identify and gracefully reject biophysically impossible or out-of-domain agroclimatic queries ($D_M(\mathbf{x}) > 5.0$).
  \item \textbf{Stochastic Multi-Objective Pareto Optimizer:} Simulates multi-year crop sequences under 100-draw Monte Carlo weather shocks, balancing economic margins, drought resilience ($\text{CVaR}_{10}$), and soil organic carbon accumulation across four actionable archetypes.
  \item \textbf{Agentic Model Context Protocol (MCP) Copilot:} Deploys five structured MCP tools streaming deterministic telemetry, live NASA Earthdata CMR granules, sub-field 3x3 DEM micro-topography grids, and external verification links without LLM math hallucination.
\end{enumerate}

\subsection{Core Research Objectives}
The core objectives guiding this investigation comprise:
\begin{enumerate}[leftmargin=*]
  \item \textbf{Physical Grounding:} Eliminate black-box hallucinations by coupling machine learning regressors directly with classical conservation equations of energy, water, and mass balance.
  \item \textbf{Uncertainty Quantification:} Formulate rotational optimization around worst-decile financial safety floors ($\text{CVaR}_{10}$) and asymmetric quantile spreads ($P10/P50/P90$) rather than fragile mean expectations.
  \item \textbf{Multi-Mission Planetary Ingestion:} Establish a resilient, low-latency data assimilation pipeline operationalizing ten spaceborne NASA missions and pedological archives without mocked data.
  \item \textbf{Global Validation \& OOD Guardrails:} Empirically benchmark algorithm performance across contrasting agroclimatic control stations (Dodoma, Fallon, Ames) while verifying Out-of-Distribution rejection against an extreme polar negative control.
\end{enumerate}

\begin{table*}[t]
\centering
\caption{Multi-Mission Satellite Remote Sensing and Global Pedotransfer Ingestion Matrix}
\label{tab:data_sources}
\small
\begin{tabularx}{\textwidth}{lcp{3.2cm}p{3.5cm}X}
\toprule
\textbf{Mission / Archive} & \textbf{Spatial Res.} & \textbf{Temporal Res.} & \textbf{Ingested Parameters} & \textbf{Agronomic Governing Function} \\
\midrule
NASA POWER v2.0 (MERRA-2) & $0.5^\circ \times 0.625^\circ$ & Daily (1981--Pres.) & $R_s$, $T_{\min}$, $T_{\max}$, $u_2$, $RH$, $P$ & FAO-56 Penman-Monteith combination evapotranspiration ($ET_0$) \\
NASA SMAP L4 (SPL4SMGP) & $9\text{ km}$ & 3-Hourly & Surface/Root Wetness ($0$--$100\text{ cm}$) & Root-zone soil bucket water balance and empirical deficit index ($WSI$) \\
NASA GPM IMERG (V07B) & $0.1^\circ \times 0.1^\circ$ & Daily & Calibrated Precipitation Flux & Kinetic rainfall erosivity factor ($R$) in RUSLE equation \\
NASA MODIS Terra (MOD13Q1) & $250\text{ m}$ & 16-Day Composite & NDVI, EVI, Surface Reflectance & Dynamic vegetative canopy coverage and bare-field exposure detection \\
NASA ECOSTRESS (L2/L4) & $70\text{ m}$ & 1--5 Days & Thermal LST, Evaporative Stress & High-resolution microclimatic thermal stress and canopy transpiration \\
NASA GEDI (L4A) & $25\text{ m}$ footprint & Mission Orbit & Aboveground Biomass Density & Structural boundary windbreak and vegetative roughness calibration \\
NASA GRACE-FO & $300\text{ km}$ & Monthly & Terrestrial Water Storage ($\Delta\text{TWS}$) & Deep aquifer groundwater depletion and regional hydraulic deficit \\
ISRIC SoilGrids 2.0 & $250\text{ m}$ & Static / Decadal & Clay, Sand, Silt, pH, Organic Carbon & Saxton-Rawls pedotransfer hydrology ($AWC, \theta_{FC}, \theta_{WP}$) \\
iSDAsoil v2 (African Basins) & $30\text{ m}$ & In-Situ Calibrated & Subsoil Texture, Micronutrients & Localized Sub-Saharan African soil parameterization and correction \\
Open-Meteo SRTM / Copernicus & $90\text{ m}$ & Static & Elevation, Slope Gradient, Aspect & Sub-field 3x3 micro-topography DEM grid and RUSLE slope factor ($LS$) \\
\bottomrule
\end{tabularx}
\end{table*}

\section{Methods}

\subsection{Multi-Mission Remote Sensing Ingestion}
TerraShift+ establishes automated, authenticated connections across eight planetary-scale Earth observation archives. Table~\ref{tab:data_sources} details the satellite missions, spatial resolutions, temporal frequencies, and biophysical utilities.

\subsection{Physics-Informed Agronomic Engine}

\subsubsection{Reference Evapotranspiration via FAO-56 Penman-Monteith}
Atmospheric evaporative demand is computed using the physically complete FAO-56 Penman-Monteith combination equation:
\begin{equation}
ET_0 = \frac{0.408 \Delta (R_n - G) + \gamma \frac{900}{T + 273} u_2 (e_s - e_a)}{\Delta + \gamma (1 + 0.34 u_2)}
\end{equation}
where $R_n$ is net radiation at the crop surface ($\text{MJ}\cdot\text{m}^{-2}\cdot\text{day}^{-1}$), $G$ is soil heat flux density (assumed zero over daily timesteps), $T$ is mean daily air temperature at 2~m ($^\circ\text{C}$), $u_2$ is wind speed at 2~m ($\text{m}\cdot\text{s}^{-1}$), $e_s$ is saturation vapor pressure ($\text{kPa}$), $e_a$ is actual vapor pressure ($\text{kPa}$), $\Delta$ is the slope of the saturation vapor pressure curve ($\text{kPa}\cdot^\circ\text{C}^{-1}$), and $\gamma$ is the psychrometric constant ($\text{kPa}\cdot^\circ\text{C}^{-1}$).

Crop evapotranspiration under non-stressed conditions ($ET_c$) and actual conditions ($ET_{c,\text{act}}$) is determined using the dual crop coefficient methodology:
\begin{align}
ET_c &= (K_{cb} + K_e) ET_0 \\
ET_{c,\text{act}} &= (K_s K_{cb} + K_e) ET_0
\end{align}
where $K_{cb}$ is the basal crop transpiration coefficient, $K_e$ is the soil evaporation coefficient, and $K_s \in [0, 1]$ is the soil water stress reduction factor.

\subsubsection{Saxton-Rawls Pedotransfer Hydrology}
Soil hydraulic properties are derived directly from soil texture (sand and clay mass fractions) and organic matter content using the Saxton-Rawls pedotransfer equations:
\begin{align}
\theta_{1500t} &= -0.024 S + 0.487 C + 0.006 OM + 0.005 (S \times OM) \nonumber \\
&\quad - 0.013 (C \times OM) + 0.068 (S \times C) + 0.031 \\
\theta_{WP} &= \theta_{1500t} + 0.14 \theta_{1500t} - 0.02 \\
\theta_{33t} &= -0.251 S + 0.195 C + 0.011 OM + 0.006 (S \times OM) \nonumber \\
&\quad - 0.027 (C \times OM) + 0.452 (S \times C) + 0.299 \\
\theta_{FC} &= \theta_{33t} + \left[1.283 (\theta_{33t})^2 - 0.374 \theta_{33t} - 0.015\right] \\
AWC &= (\theta_{FC} - \theta_{WP}) \times Z_r
\end{align}
where $S$ is sand fraction, $C$ is clay fraction, $OM$ is organic matter fraction, $\theta_{FC}$ is volumetric water content at field capacity ($-33\text{ kPa}$), $\theta_{WP}$ is wilting point ($-1500\text{ kPa}$), and $Z_r$ is dynamic crop rooting depth ($mm$).

\subsubsection{Dynamic Root-Zone Water Balance \& Stress Index}
Daily root-zone depletion $D_r$ evolves according to the mass balance equation:
\begin{equation}
D_{r,t} = D_{r,t-1} - (P_t - RO_t) - I_t - CR_t + ET_{c,\text{act},t} + DP_t
\end{equation}
where $P_t$ is precipitation, $RO_t$ is surface runoff, $I_t$ is irrigation, $CR_t$ is capillary rise from groundwater, and $DP_t$ is deep percolation. The daily Crop Water Stress Index ($WSI$) is defined as:
\begin{equation}
WSI_t = 1.0 - \frac{ET_{c,\text{act},t}}{ET_{c,t}}
\end{equation}
Integrating $WSI_t$ over the growth cycle weighted by phenological stage sensitivity factors ($k_y$) yields cumulative yield reduction penalties.

\subsubsection{Thermal Phenology: GDD and KDD}
Vegetative development and heat injury are governed by thermal integration:
\begin{align}
GDD_t &= \max\left(0, \frac{T_{\max,t} + T_{\min,t}}{2} - T_{\text{base}}\right) \\
KDD_t &= \max\left(0, T_{\max,t} - T_{\text{kill}}\right)
\end{align}
where $T_{\text{base}}$ represents crop-specific minimum developmental threshold ($10^\circ\text{C}$ for maize, $8^\circ\text{C}$ for sorghum) and $T_{\text{kill}}$ represents the critical thermal denaturation threshold ($35^\circ\text{C}$).

\subsubsection{Topographic Soil Erosion via RUSLE}
Mean annual soil loss $A$ ($\text{t}\cdot\text{ha}^{-1}\cdot\text{yr}^{-1}$) is modeled through the Revised Universal Soil Loss Equation:
\begin{equation}
A = R \times K \times LS \times C \times P
\end{equation}
where $R$ is rainfall erosivity calculated from GPM IMERG, $K$ is soil erodibility from ISRIC SoilGrids, $LS$ is the slope length and steepness factor computed from the DEM, $C$ is the cover-management factor dynamically assigned based on crop canopy coverage, and $P$ is the erosion control practice factor. Under bare fallow, $C \to 1.0$, producing severe soil detachment; under diversified cover cropping, $\mathrm{RUSLE}_{\mathrm{cover}} < \mathrm{RUSLE}_{\mathrm{bare}}$, cutting erosion rates by up to 65\%.

\subsection{Machine Learning Surrogate \& Control Validation}

\subsubsection{Quantile Gradient Boosting Yield Surrogates}
To deliver instantaneous rotational evaluations while bounding predictive uncertainty, TerraShift+ deploys Quantile Gradient Boosting Regressors trained with the asymmetric pinball loss:
\begin{equation}
\mathcal{L}_\alpha(y, \hat{y}) = \max\left(\alpha (y - \hat{y}), (1 - \alpha)(\hat{y} - y)\right)
\end{equation}
Three separate surrogate ensembles are trained for $\alpha \in \{0.10, 0.50, 0.90\}$, providing lower bound ($P10$), median expected ($P50$), and optimistic ($P90$) yield and gross margin estimates.

\subsubsection{GMM Mahalanobis Out-of-Distribution Guardrail}
To prevent ungrounded predictions when presented with non-agricultural or geographically corrupt queries, a 72-component Gaussian Mixture Model (GMM) is fitted over the global training feature manifold. The Mahalanobis distance $D_M(\mathbf{x})$ to the nearest cluster $k$:
\begin{equation}
D_M(\mathbf{x}) = \min_k \sqrt{(\mathbf{x} - \mu_k)^T \Sigma_k^{-1} (\mathbf{x} - \mu_k)}
\end{equation}
If $D_M(\mathbf{x}) > \tau_{\text{OOD}} = 5.0$, the input is flagged as Out-of-Distribution, preventing black-box extrapolation and triggering deterministic physical fallback diagnostics.

\subsection{Stochastic Multi-Objective Pareto Optimizer}
The rotation planner evaluates all permissible 5-year crop sequences $\mathbf{s} = (c_1, c_2, \dots, c_5)$ across a candidate crop library $\mathcal{C}$. To enforce rotational agronomy:
\begin{equation}
c_{t} \neq c_{t-1} \quad \forall t \in \{2,\dots,5\} \quad \text{if } c_{t} \in \mathcal{C}_{\text{depleting}}
\end{equation}
Continuous monoculture produces an acute nitrogen penalty ($\mathrm{Net}\,N_{\mathrm{mono}} < 0$). In contrast, leguminous species ($c \in \mathcal{C}_{\text{legume}}$) contribute a biological nitrogen credit ($\Delta N = +45\text{ kg}\cdot\text{ha}^{-1}$) that reduces synthetic fertilizer requirements for the subsequent crop cycle.

Multi-objective fitness evaluates three competing dimensions under $N=100$ Monte Carlo climate shock draws:
\begin{align}
f_1(\mathbf{s}) &= \mathbb{E}[\text{Margin}_{5Y}(\mathbf{s})] = \frac{1}{N}\sum_{i=1}^N \sum_{t=1}^5 M(c_t, \omega_{i,t}) \\
f_2(\mathbf{s}) &= \text{CVaR}_{10}(\mathbf{s}) = \mathbb{E}\left[\text{Margin}_{5Y}(\mathbf{s}) \;\middle|\; \text{Margin}_{5Y} \le q_{0.10}\right] \\
f_3(\mathbf{s}) &= \sum_{t=1}^5 \left( \Delta SOC(c_t) - \frac{A(c_t)}{20} \right)
\end{align}
The solver constructs the non-dominated Pareto frontier $\mathcal{P}^*$ and extracts four operational archetypes:
\begin{itemize}[leftmargin=*]
  \item \textbf{Commercial Maximum Profit:} Maximizes $f_1(\mathbf{s})$ subject to $WSI \le 0.45$.
  \item \textbf{Climate Adaptation \& Drought-Resilient:} Maximizes $f_2(\mathbf{s})$ ($\text{CVaR}_{10}$) under acute water deficit conditions.
  \item \textbf{Soil Restorative \& Conservation:} Maximizes $f_3(\mathbf{s})$ to enhance topsoil organic matter and minimize RUSLE detachment.
  \item \textbf{Balanced Sustainable:} Computes the centroid compromise vector across the Pareto front.
\end{itemize}

\subsection{Agentic Model Context Protocol (MCP) Copilot}
TerraShift+ includes an agentic copilot executing five Model Context Protocol (MCP) tools:
\begin{enumerate}[leftmargin=*]
  \item \texttt{nasa\_power\_agroclimatology(lat, lon)}: Fetches daily climatology, 5-year trend lines, and raw NASA POWER JSON API links.
  \item \texttt{nasa\_cmr\_and\_worldview\_lookup(lat, lon, layer)}: Queries Earthdata CMR for active mission granules and constructs localized NASA Worldview deep links.
  \item \texttt{subfield\_dem\_vulnerability\_grid(lat, lon, metric)}: Derives the 3x3 DEM micro-topography grid, calculating slope, runoff vectors, and water stress hotspots.
  \item \texttt{soil\_and\_crop\_profile(lat, lon)}: Extracts soil texture fractions, pH, and organic carbon from iSDAsoil/SoilGrids.
  \item \texttt{search\_nasa\_earthdata\_knowledge(query)}: Semantic FTS5 retrieval over 45 peer-reviewed NASA Earth science collections and DOIs.
\end{enumerate}

\section{Results}

\subsection{Spatial Cross-Validation Performance}
The quantile regression surrogate models were validated using 10-fold Spatial Block Cross-Validation across 72 global agricultural clusters to prevent spatial autocorrelation leakage. Performance metrics are reported in Table~\ref{tab:surrogate_performance}.

\begin{table}[htbp]
\centering
\caption{Spatial Cross-Validation Performance Across 72 Agroecological Clusters}
\label{tab:surrogate_performance}
\small
\begin{tabular}{lcccc}
\toprule
\textbf{Target Metric} & \textbf{RMSE} & \textbf{$R^2$} & \textbf{Pinball Loss} & \textbf{Coverage} \\
\midrule
Yield Median ($P50$) & $0.050$ & $0.942$ & --- & --- \\
Yield Lower ($P10$)  & --- & --- & $0.009$ & $9.8\%$ \\
Yield Upper ($P90$)  & --- & --- & $0.008$ & $89.2\%$ \\
Empirical $80\%$ Interval & --- & --- & --- & $83.5\%$ \\
Water Stress ($WSI$) & $0.041$ & $0.958$ & $0.007$ & $82.1\%$ \\
RUSLE Soil Loss ($A$) & $1.12\text{ t/ha}$ & $0.916$ & $0.024$ & $84.0\%$ \\
\bottomrule
\end{tabular}
\end{table}

\subsection{Benchmark Control Station Validation}
The pipeline was evaluated across four canonical stations:
\begin{enumerate}[leftmargin=*]
  \item \textbf{Dodoma, Tanzania ($-6.1630, 35.7516$):} Tropical semiarid savanna experiencing severe unimodal dry spells.
  \item \textbf{Fallon, Nevada ($39.4735, -118.7774$):} Arid Great Basin high-desert irrigated agriculture.
  \item \textbf{Ames, Iowa ($42.0308, -93.6319$):} Humid continental high-yield maize-soybean corn belt.
  \item \textbf{Polar Negative Control ($78.2232, 15.6267$):} Svalbard Arctic permafrost testing Out-of-Distribution rejection.
\end{enumerate}

\begin{table*}[t]
\centering
\caption{Empirical Benchmark Across Canonical Validation Stations}
\label{tab:benchmark_stations}
\small
\begin{tabularx}{\textwidth}{lXXXX}
\toprule
\textbf{Metric / Parameter} & \textbf{Dodoma, Tanzania} & \textbf{Fallon, Nevada} & \textbf{Ames, Iowa} & \textbf{Polar Control (Svalbard)} \\
\midrule
Agroclimatic Zone & Semiarid Tropical Savanna & Arid Cold Desert Basin & Humid Continental Corn Belt & Arctic High-Latitude Tundra \\
Annual Precipitation & $831.2\text{ mm}$ & $128.4\text{ mm}$ & $924.5\text{ mm}$ & $210.3\text{ mm}$ \\
Annual FAO-56 $ET_0$ & $2500.2\text{ mm}$ & $1420.5\text{ mm}$ & $1050.2\text{ mm}$ & $320.1\text{ mm}$ \\
Precipitation Deficit ($P - ET_0$) & $-1669.0\text{ mm}$ & $-1292.1\text{ mm}$ & $-125.7\text{ mm}$ & $-109.8\text{ mm}$ \\
SMAP L4 Root Wetness ($Z$) & $-0.23$ (Deficit) & $-0.85$ (Extreme Arid) & $+0.32$ (Adequate) & N/A (Frozen Permafrost) \\
Growing Degree Days ($GDD$) & $3420\text{ }^\circ\text{C}\cdot\text{d}$ & $1840\text{ }^\circ\text{C}\cdot\text{d}$ & $1680\text{ }^\circ\text{C}\cdot\text{d}$ & $12\text{ }^\circ\text{C}\cdot\text{d}$ \\
GMM Mahalanobis Distance & $1.15$ (In-Domain) & $2.42$ (In-Domain) & $0.85$ (In-Domain) & \textbf{8.40 (OOD Rejected)} \\
Recommended Archetype & Climate Adaptation & Soil Restorative / Forage & Commercial Maximum & Rejection Diagnostic \\
5-Year Crop Sequence & Pigeonpea-Millet-Sorghum & Alfalfa-Teff-Forage & Maize-Soybean-Maize & None (Below GDD Threshold) \\
Baseline Monoculture WSI & $0.288$ & $0.580$ & $0.140$ & N/A \\
Optimal Rotation WSI & $0.222$ ($-22.9\%$) & $0.340$ ($-41.4\%$) & $0.110$ ($-21.4\%$) & N/A \\
5-Year Expected Margin & $\$2,480\text{/ha}$ & $\$1,850\text{/ha}$ & $\$4,120\text{/ha}$ & N/A \\
Downside Risk ($\text{CVaR}_{10}$) & $\$1,720\text{/ha}$ & $\$1,100\text{/ha}$ & $\$3,250\text{/ha}$ & N/A \\
RUSLE Soil Loss Reduction & $-64.5\%$ & $-58.2\%$ & $-42.0\%$ & N/A \\
\bottomrule
\end{tabularx}
\end{table*}

\subsection{Drought Shock Monotonicity Stress-Testing}
To verify thermodynamic and biological consistency, the optimizer was subjected to synthetic climate stress tests at Dodoma, Tanzania, progressively degrading annual precipitation by $-25\%$ while elevating daily mean temperatures by $+2.0^\circ\text{C}$.

\begin{table}[htbp]
\centering
\caption{Monotonicity Response Under Simulated Drought Shocks (Dodoma)}
\label{tab:drought_shock}
\small
\begin{tabular}{lcccc}
\toprule
\textbf{Climate Shock} & \textbf{Precip.} & \textbf{Mean WSI} & \textbf{Margin ($P50$)} & \textbf{CVaR$_{10}$} \\
\midrule
Baseline Climatology & $831.2\text{ mm}$ & $0.222$ & $\$2,480$ & $\$1,720$ \\
Mild Shock ($-10\%$) & $748.1\text{ mm}$ & $0.265$ & $\$2,210$ & $\$1,450$ \\
Severe Shock ($-25\%$) & $623.4\text{ mm}$ & $0.341$ & $\$1,840$ & $\$980$ \\
Compound ($-25\%, +2^\circ\text{C}$) & $623.4\text{ mm}$ & $0.412$ & $\$1,520$ & $\$650$ \\
\bottomrule
\end{tabular}
\end{table}

As shown in Table~\ref{tab:drought_shock}, water stress strictly increases monotonically ($\partial WSI / \partial P \le 0$), while financial returns and risk bounds decrease predictably without non-physical oscillations.

\subsection{Latency \& Throughput Benchmarking}
Real-time operational suitability was benchmarked across 1,000 requests. Cold-start requests (fetching fresh NASA POWER/SMAP telemetry via HTTP) resolve in $585.2\text{ ms}$, while cached in-memory evaluations resolve in $1.50\text{ ms}$ ($>660\text{ req/sec}$), fully enabling interactive web and mobile deployment.

\section{Discussion}

\subsection{Scientific Significance}
TerraShift+ demonstrates that embedding biophysical conservation laws into machine learning pipelines resolves the trade-off between the adaptability of empirical data and the physical rigor of deterministic models. By bounding surrogate neural and boosted tree predictions within FAO-56 Penman-Monteith, Saxton-Rawls, and RUSLE governing limits, the system completely prevents agronomic hallucinations while sustaining sub-second execution speeds.

\begin{table*}[t]
\centering
\caption{System Comparison: Static Almanacs vs. Black-Box LLMs vs. TerraShift+}
\label{tab:system_comparison}
\small
\begin{tabularx}{\textwidth}{lXXX}
\toprule
\textbf{Capability / Dimension} & \textbf{Static Crop Almanacs} & \textbf{Black-Box Generative LLMs} & \textbf{TerraShift+ Closed-Loop PIML} \\
\midrule
Climate Non-Stationarity & Complete failure under extremes & Plausible prose, uncalibrated math & Dynamic 5-year Theil-Sen drift telemetry \\
Mass / Water Conservation & Ignored & Frequently violated ($WSI < 0$ or $>1$) & Strictly bounded via FAO-56 \& Saxton-Rawls \\
Multi-Objective Optimization & Single-crop heuristics & Unranked text suggestions & Formal 4-archetype Pareto frontier \\
Uncertainty Quantification & None & None (deterministic hallucinations) & Quantile regression ($P10/P50/P90$) and $\text{CVaR}_{10}$ \\
Out-of-Distribution Safety & None & High hallucination risk in novel domains & GMM Mahalanobis distance rejection ($D_M > 5.0$) \\
Bilingual African Dialects & English / Colonial only & Generic Swahili without agro-lexicon & Verified Tanzanian Swahili (`sw-TZ`) \\
Telemetry Provenance & None & Uncited or fictional web links & Direct NASA Earthdata CMR \& Worldview DOIs \\
\bottomrule
\end{tabularx}
\end{table*}

\subsection{Comparison with Existing Paradigms}
Table~\ref{tab:system_comparison} contrasts TerraShift+ against static almanacs and black-box generative models. TerraShift+ uniquely bridges planetary remote sensing and field-scale operational decision support.

\subsection{Limitations and Future Work}
Ingested global soil databases (ISRIC SoilGrids and iSDAsoil) carry local uncertainty margins of up to $\pm15\%$ prior to field-level calibration. Ongoing efforts focus on integrating mobile smartphone-based spectroscopy and LoRaWAN soil moisture probes to further ground predictive models.

\section{Conclusions}
TerraShift+ establishes a robust, physics-informed remote sensing architecture that empowers smallholder and commercial farmers to navigate accelerating climatic volatility. By fusing ten NASA satellite missions with rigorous physical laws, quantile surrogate regression, and multi-objective Pareto optimization, the platform delivers verifiable, resilient crop rotation sequences that protect topsoil, conserve scarce water resources, and secure agricultural livelihoods worldwide.

\section*{NASA Space Apps Challenge AI Provenance Disclosure}
In compliance with NASA Space Apps Challenge guidelines: Large Language Models (Anthropic Claude, Google Gemini) were utilized within strictly bounded retrieval-augmented generation (RAG) guardrails. LLMs were prohibited from synthesizing quantitative yield, water stress, or financial metrics independently; all numerical outputs originated exclusively from deterministic biophysical equations and validated surrogate models. AI coding assistants were utilized for unit test generation and LaTeX formatting under active human engineering supervision.

\begin{thebibliography}{15}
\small
\bibitem{allen1998}
R.~G. Allen, L.~S. Pereira, D.~Raes, and M.~Smith, ``Crop evapotranspiration: Guidelines for computing crop water requirements,'' \textit{FAO Irrigation and Drainage Paper 56}, Rome, Italy, 1998.

\bibitem{saxton2006}
K.~E. Saxton and W.~J. Rawls, ``Soil water characteristic estimates by texture and organic matter for hydrologic solutions,'' \textit{Soil Sci. Soc. Am. J.}, vol.~70, no.~5, pp.~1569--1578, 2006.

\bibitem{renard1997}
K.~G. Renard, G.~R. Foster, G.~A. Weesies, D.~K. McCool, and D.~C. Yoder, ``Predicting soil erosion by water: A guide to conservation planning with the Revised Universal Soil Loss Equation (RUSLE),'' \textit{USDA Agricultural Handbook No. 703}, 1997.

\bibitem{power2024}
P.~W. Stackhouse et~al., ``NASA POWER Release 2.0: Daily agroclimatology methodology,'' NASA Langley Research Center, Hampton, VA, 2024. [Online]. Available: \url{https://power.larc.nasa.gov/}

\bibitem{smap2024}
R.~H. Reichle et~al., ``SMAP L4 Global 9 km EASE-Grid Surface and Root-Zone Soil Moisture, Version 8,'' NASA National Snow and Ice Data Center (NSIDC) DAAC, Boulder, CO, 2024. \url{https://doi.org/10.5067/T5RUATAQREF8}

\bibitem{ecostress2024}
G.~Hulley et~al., ``ECOSTRESS Land Surface Temperature and Emissivity Daily L2 Global 70 m,'' NASA LP DAAC, 2024. \url{https://doi.org/10.5067/ECOSTRESS/ECO2LSTE.001}

\bibitem{gpm2024}
G.~J. Huffman et~al., ``GPM IMERG Final Precipitation L3 1 month 0.1 degree x 0.1 degree V07,'' NASA GES DISC, 2024. \url{https://doi.org/10.5067/GPM/IMERG/3B-MONTH/07}

\bibitem{modis2024}
K.~Didan, ``MOD13Q1 MODIS/Terra Vegetation Indices 16-Day L3 Global 250 m SIN Grid V061,'' NASA LP DAAC, 2024. \url{https://doi.org/10.5067/MODIS/MOD13Q1.061}

\bibitem{gedi2024}
R.~Dubayah et~al., ``GEDI L4A Footprint Aboveground Biomass Density, Version 2.1,'' NASA ORNL DAAC, 2024. \url{https://doi.org/10.3334/ORNLDAAC/2056}

\bibitem{isdasoil2021}
M.~G. Walsh et~al., ``High-resolution priority soil property maps for Sub-Saharan Africa (iSDAsoil),'' \textit{Sci. Data}, vol.~8, p.~255, 2021. \url{https://doi.org/10.1038/s41597-021-01033-6}

\bibitem{soilgrids2020}
L.~Poggio et~al., ``SoilGrids 2.0: producing soil information for the globe with quantified spatial uncertainty,'' \textit{SOIL}, vol.~7, pp.~217--240, 2021. \url{https://doi.org/10.5194/soil-7-217-2021}

\bibitem{koenker2005}
R.~Koenker, \textit{Quantile Regression}. Cambridge University Press, Cambridge, UK, 2005.

\bibitem{rockafellar2000}
R.~T. Rockafellar and S.~Uryasev, ``Optimization of conditional value-at-risk,'' \textit{J. Risk}, vol.~2, no.~3, pp.~21--41, 2000.

\bibitem{karniadakis2021}
G.~E. Karniadakis et~al., ``Physics-informed machine learning,'' \textit{Nat. Rev. Phys.}, vol.~3, pp.~422--440, 2021.
\end{thebibliography}

\appendix
\section{Data Availability \& Verified Endpoints}
All satellite datasets, REST endpoints, and model parameter files are publicly accessible without authentication barriers:
\begin{itemize}[leftmargin=*]
  \item NASA POWER: \url{https://power.larc.nasa.gov/api/temporal/daily/point}
  \item NSIDC SMAP L4 DOI: \url{https://doi.org/10.5067/T5RUATAQREF8}
  \item NASA Worldview: \url{https://worldview.earthdata.nasa.gov/}
  \item ISRIC SoilGrids 2.0: \url{https://rest.isric.org/soilgrids/v2.0/}
  \item Open-Meteo DEM 90m: \url{https://api.open-meteo.com/v1/elevation}
  \item TerraShift+ Live App: \url{https://terrashift-plus.pages.dev}
  \item Interactive Swagger Docs: \url{https://terrashift-plus.pages.dev/docs.html}
\end{itemize}

\section{System Requirements \& Edge Deployment}
\subsection{Cloudflare Pages Edge Runtime (Production)}
\begin{itemize}[leftmargin=*]
  \item Runtime: Cloudflare Workers V8 JavaScript Engine
  \item CPU Budget: $<10\text{ ms}$ per invocation
  \item Memory Budget: $128\text{ MB}$
  \item Cold-Start Latency: $<600\text{ ms}$ (including upstream NASA fetch)
  \item In-Memory Cache Latency: $1.50\text{ ms}$ ($>660\text{ req/sec}$)
\end{itemize}

\subsection{Local Development Core}
\begin{itemize}[leftmargin=*]
  \item Python 3.11+ (FastAPI, NumPy, SciPy, scikit-learn, Pydantic)
  \item Node.js 18+ (React 18, Vite, Tailwind CSS, Leaflet)
\end{itemize}

\section{Open-Source Software Dependencies}
The platform builds exclusively on standard open-source libraries:
\begin{itemize}[leftmargin=*]
  \item \textbf{Core Modeling:} \texttt{numpy} (v1.26.4), \texttt{scipy} (v1.12.0), \texttt{scikit-learn} (v1.4.1), \texttt{pandas} (v2.2.1)
  \item \textbf{Web Engine:} \texttt{fastapi} (v0.110.0), \texttt{uvicorn} (v0.28.0), \texttt{pydantic} (v2.6.4)
  \item \textbf{Frontend UI:} \texttt{react} (v18.3.1), \texttt{vite} (v6.4.6), \texttt{leaflet} (v1.9.4), \texttt{lucide-react} (v0.344.0)
\end{itemize}

\section{Numerical Reproducibility Checklist}
\begin{itemize}[leftmargin=*]
  \item Random seed fixed to $42$ across all Monte Carlo draws and Spatial Block Cross-Validation splits.
  \item Standardized model weights exported to \texttt{functions/\_lib/model\_weights.json}.
  \item Control station validation reproducible via: \\
  \texttt{python backend/training/run\_paper\_validation.py}
  \item Full test suite passes: \\
  \texttt{pytest backend/tests/test\_nextgen\_chat.py}
\end{itemize}

\end{document}
